% !TEX root = ../thesis_book_0421052099.tex Every result quoted here comes from
% results/phase4/thesis_numbers.json, the same source the body chapters read
% from. The two exceptions are the entity and triple counts, which are not in
% that file: they come from results/phase1/ by way of tab:graphstats. If a rerun
% changes a value in either source, this file is one of the places that has to
% change with it.

Large language models answer factual questions fluently. But, they tend to sound
just as fluent when they don't hold the facts. To fix this, researchers
introduced Retrieval-Augmented Generation (RAG), which connects the LLM to an
external data source. This grounds the model's answers in retrieved text. But,
the retrieval happens once and up front, before any reasoning starts. So, a
question that spans several facts has to be answered from whatever that first
query brought back. Systems that navigate a knowledge graph do better by
interleaving retrieval with reasoning. Even they, however, send out whatever
their final generation call produces, and nothing checks the answer's claims
before it reaches the user.

To close that gap we propose Agentic Graph Reasoning (AGR): a knowledge-graph
question-answering framework. Here, an agent plans sub-objectives and walks the
graph through a constrained, deterministic tool API. Its distinguishing
component is a Structural Verification Layer. Before the answer is emitted, it
splits the draft into atomic claims and checks each one against the (head,
relation, tail) triples the agent actually traversed. Claims it cannot ground
are re-explored or withdrawn. So, where it cannot ground a claim, AGR hedges
rather than asserts. And the grounded claims come paired with the triples that
support them.

We evaluate AGR on a Freebase-derived environment of $2.59$ million entities and
$8.31$ million triples. We test it on $400$ certified questions from each of
WebQSP and ComplexWebQuestions. We compare it against four baselines: parametric
answering, vector retrieval, GraphRAG, and an agentic navigator. All five
systems share one frozen backbone and the same call budget. So, any difference
between them comes from architecture, not model capacity.

AGR reaches a Hits@1 of $0.755$ on WebQSP and $0.522$ on ComplexWebQuestions. It
beats every baseline on both datasets, at half the agentic navigator's
language-model calls. The navigator does lead AGR on the questions its own
budget lets it finish. But that budget runs out on $29\%$ and $44\%$ of
questions on the two datasets. AGR's own budget never does. AGR also grounds
every entity it asserts, all $1{,}709$ of them. The parametric control leaves
$221$ of its $1{,}001$ asserted entities ungrounded ($22.1\%$). And on
ComplexWebQuestions, AGR is the only system whose accuracy rises with hop count.
Every other system ends up below where it started.

We performed an ablation study. Each condition removes exactly one component and
re-runs the identical questions. The four components are the planner that
decomposes a question into sub-objectives, the backtracker, the claim
verification layer, and the model-based edge scorer. We tested the re-runs
against the unmodified system with paired significance tests. Only the planner
shows a detectable effect, and it runs the opposite way from what we expected.
Removing it \emph{improves} WebQSP accuracy by $0.083$ F1 ($p = 0.006$) and cuts
token use by $31\%$, while it trends the other way on ComplexWebQuestions. So,
decomposition should be gated on question structure, not applied
unconditionally.

We read all $256$ remaining failures by hand. The characteristic error is not
invention. We call this the \emph{echo attractor}. It is a real, grounded entity
one hop from the answer, which any grounding check will pass because it is
genuinely true, and which different systems fall into together. The same reading
confirms $57$ questions where the benchmark, not the system, was wrong.

\endinput
